% GENERATED FILE. Edit canonical lessons, then run npm run generate:books.
% canonical-source-hash: fnv1a64:aa735594e5122ae5
% canonical-lessons: GU-C370-khatro, GU-C370-hani, GU-C370-apatti, GU-C370-bagadvu-spoil, GU-C370-vethvu, GU-R370-review-363-370

\chapter{Danger and Harm}
\label{ch:gu-danger-and-harm}
\hlchaptermodality{370}

\begin{chapteropening}
\textbf{By the end of this chapter:} \emph{I can say what harm will follow if nothing is done.}

Complete the last lesson of chapter 370: I can say what harm will follow if nothing is done.
\end{chapteropening}

\section[\texorpdfstring{\gu{ખતરો} (khatro)}{ખતરો (khatro)}]{\gu{ખતરો} (khatro) — a peril}

\label{lesson:GU-C370-khatro}



\begin{quote}
Before the new one: say the Gujarati for a misgiving, then the Gujarati for to frighten.
\end{quote}

\subsection*{You'll want to know: \gu{ખતરો}}
\textbf{\gu{ખતરો}} — \emph{khatro} — \textquotedblleft{}a peril\textquotedblright{}.

\begin{quote}
\gu{જો} \gu{આ} \gu{જૂના} \gu{પુલ} \gu{પર} \gu{ભારે} \gu{ગાડીઓ} \gu{ચાલશે} \gu{તો} \gu{ખતરો} \gu{છે}.
\end{quote}

\begin{quote}
\emph{jo ā jūnā pul par bhāre gāḍīo cālśe to khatro che.}
\end{quote}

\begin{quote}
If heavy vehicles use this old bridge, it is dangerous.
\end{quote}

A masculine noun.

\begin{cousinweb}
\textbf{\gu{ખતરો}} (\emph{khatro}) is usually traced to Arabic \emph{khaṭar}, \textquotedblleft{}danger, risk\textquotedblright{}, through Persian.
\end{cousinweb}

\begin{grammarlens}[title={What you've built}]
One more word for saying what happens if.
\end{grammarlens}

\subsection*{Your turn}
\begin{itemize}
\raggedright
  \item \emph{Say it:} \emph{khatro}
  \item \emph{Say it:} \emph{khatro}, once more
  \item \emph{Say it:} \emph{jo ā jūnā pul par bhāre gāḍīo cālśe to khatro che.} — If heavy vehicles use this old bridge, it is dangerous.
  \item \emph{Recall:} say the Gujarati for a misgiving, then the Gujarati for to frighten, then say \emph{khatro} again
\end{itemize}

\begin{tcolorbox}[breakable,colback=teal!4,colframe=teal!35!black,title={Before you move on}]
What does \gu{ખતરો} mean? (\textquotedblleft{}A peril\textquotedblright{}.) Say it once more.
\end{tcolorbox}
\section[\texorpdfstring{\gu{હાનિ} (hāni)}{હાનિ (hāni)}]{\gu{હાનિ} (hāni) — harm}

\label{lesson:GU-C370-hani}



\begin{quote}
Before the new one: say the Gujarati for to frighten, then the Gujarati for a peril.
\end{quote}

\subsection*{You'll want to know: \gu{હાનિ}}
\textbf{\gu{હાનિ}} — \emph{hāni} — \textquotedblleft{}harm\textquotedblright{}.

\begin{quote}
\gu{જો} \gu{તમે} \gu{વધારે} \gu{દવા} \gu{લેશો} \gu{તો} \gu{શરીરને} \gu{હાનિ} \gu{થશે}.
\end{quote}

\begin{quote}
\emph{jo tame vadhāre davā leśo to śarīrne hāni thaśe.}
\end{quote}

\begin{quote}
If you take too much medicine, it will harm your body.
\end{quote}

A feminine noun.

\begin{cousinweb}
\textbf{\gu{હાનિ}} (\emph{hāni}) is a tatsama, borrowed from Sanskrit \emph{hāni}, \textquotedblleft{}abandonment, loss, damage\textquotedblright{} (Monier-Williams).
\end{cousinweb}

\begin{grammarlens}[title={What you've built}]
One more word for saying what happens if.
\end{grammarlens}

\subsection*{Your turn}
\begin{itemize}
\raggedright
  \item \emph{Say it:} \emph{hāni}
  \item \emph{Say it:} \emph{hāni}, once more
  \item \emph{Say it:} \emph{jo tame vadhāre davā leśo to śarīrne hāni thaśe.} — If you take too much medicine, it will harm your body.
  \item \emph{Recall:} say the Gujarati for to frighten, then the Gujarati for a peril, then say \emph{hāni} again
\end{itemize}

\begin{tcolorbox}[breakable,colback=teal!4,colframe=teal!35!black,title={Before you move on}]
What does \gu{હાનિ} mean? (\textquotedblleft{}Harm\textquotedblright{}.) Say it once more.
\end{tcolorbox}
\section[\texorpdfstring{\gu{આપત્તિ} (āpatti)}{આપત્તિ (āpatti)}]{\gu{આપત્તિ} (āpatti) — a misfortune; a calamity}

\label{lesson:GU-C370-apatti}



\begin{quote}
Before the new one: say the Gujarati for a peril, then the Gujarati for harm.
\end{quote}

\subsection*{You'll want to know: \gu{આપત્તિ}}
\textbf{\gu{આપત્તિ}} — \emph{āpatti} — \textquotedblleft{}a misfortune; a calamity\textquotedblright{}.

\begin{quote}
\gu{જો} \gu{આ} \gu{વરસે} \gu{પણ} \gu{દુકાળ} \gu{પડશે} \gu{તો} \gu{ગામ} \gu{પર} \gu{મોટી} \gu{આપત્તિ} \gu{આવશે}.
\end{quote}

\begin{quote}
\emph{jo ā varse paṇ dukāḷ paḍśe to gām par moṭī āpatti āvśe.}
\end{quote}

\begin{quote}
If there is a drought again this year, a great calamity will strike the village.
\end{quote}

A feminine noun.

\begin{cousinweb}
\textbf{\gu{આપત્તિ}} (\emph{āpatti}) is a tatsama, borrowed from Sanskrit \emph{āpatti}, \textquotedblleft{}happening, occurring; misfortune\textquotedblright{} (Monier-Williams).
\end{cousinweb}

\begin{grammarlens}[title={What you've built}]
One more word for saying what happens if.
\end{grammarlens}

\subsection*{Your turn}
\begin{itemize}
\raggedright
  \item \emph{Say it:} \emph{āpatti}
  \item \emph{Say it:} \emph{āpatti}, once more
  \item \emph{Say it:} \emph{jo ā varse paṇ dukāḷ paḍśe to gām par moṭī āpatti āvśe.} — If there is a drought again this year, a great calamity will strike the village.
  \item \emph{Recall:} say the Gujarati for a peril, then the Gujarati for harm, then say \emph{āpatti} again
\end{itemize}

\begin{tcolorbox}[breakable,colback=teal!4,colframe=teal!35!black,title={Before you move on}]
What does \gu{આપત્તિ} mean? (\textquotedblleft{}A misfortune; a calamity\textquotedblright{}.) Say it once more.
\end{tcolorbox}
\section[\texorpdfstring{\gu{બગાડવું} (bagāḍvũ)}{બગાડવું (bagāḍvũ)}]{\gu{બગાડવું} (bagāḍvũ) — to spoil}

\label{lesson:GU-C370-bagadvu-spoil}



\begin{quote}
Before the new one: say the Gujarati for harm, then the Gujarati for a misfortune.
\end{quote}

\subsection*{You'll want to know: \gu{બગાડવું}}
\textbf{\gu{બગાડવું}} — \emph{bagāḍvũ} — \textquotedblleft{}to spoil\textquotedblright{}.

\begin{quote}
\gu{જો} \gu{તમે} \gu{વધારે} \gu{મીઠું} \gu{નાખશો} \gu{તો} \gu{દાળ} \gu{બગાડશો}.
\end{quote}

\begin{quote}
\emph{jo tame vadhāre mīṭhũ nākhśo to dāḷ bagāḍśo.}
\end{quote}

\begin{quote}
If you put in too much salt, you will spoil the dal.
\end{quote}

Past, with \gu{તેણે} (he or she …): \gu{તેણે} \gu{આખું} \gu{કામ} \gu{બગાડ્યું}, \emph{teṇe ākhũ kām bagāḍyũ}.

\begin{cousinweb}
\textbf{\gu{બગાડવું}} (\emph{bagāḍvũ}) is the transitive partner of \textbf{\gu{બગડવું}} (\emph{bagaḍvũ}), \textquotedblleft{}to go bad\textquotedblright{}, which CDIAL traces to Sanskrit \emph{vighaṭatē}, \textquotedblleft{}flies apart\textquotedblright{} (CDIAL 11673).
\end{cousinweb}

\begin{grammarlens}[title={What you've built}]
One more word for saying what happens if.
\end{grammarlens}

\subsection*{Your turn}
\begin{itemize}
\raggedright
  \item \emph{Say it:} \emph{bagāḍvũ}
  \item \emph{Say it:} \emph{bagāḍvũ}, once more
  \item \emph{Say it:} \emph{jo tame vadhāre mīṭhũ nākhśo to dāḷ bagāḍśo.} — If you put in too much salt, you will spoil the dal.
  \item \emph{Recall:} say the Gujarati for harm, then the Gujarati for a misfortune, then say \emph{bagāḍvũ} again
\end{itemize}

\begin{tcolorbox}[breakable,colback=teal!4,colframe=teal!35!black,title={Before you move on}]
What does \gu{બગાડવું} mean? (\textquotedblleft{}To spoil\textquotedblright{}.) Say it once more.
\end{tcolorbox}
\section[\texorpdfstring{\gu{વેઠવું} (veṭhvũ)}{વેઠવું (veṭhvũ)}]{\gu{વેઠવું} (veṭhvũ) — to put up with (hardship)}

\label{lesson:GU-C370-vethvu}



\begin{quote}
Before the new one: say the Gujarati for a misfortune, then the Gujarati for to spoil.
\end{quote}

\subsection*{You'll want to know: \gu{વેઠવું}}
\textbf{\gu{વેઠવું}} — \emph{veṭhvũ} — \textquotedblleft{}to put up with (hardship)\textquotedblright{}.

\begin{quote}
\gu{જો} \gu{તમે} \gu{અત્યારે} \gu{મહેનત} \gu{કરશો} \gu{તો} \gu{પછી} \gu{દુઃખ} \gu{વેઠવું} \gu{નહીં} \gu{પડે}.
\end{quote}

\begin{quote}
\emph{jo tame atyāre mahenat karśo to pachī duḥkh veṭhvũ nahī̃ paḍe.}
\end{quote}

\begin{quote}
If you work hard now, you will not have to suffer later.
\end{quote}

Past, with \gu{તેણે} (he or she …): \gu{તેણે} \gu{ઘણું} \gu{દુઃખ} \gu{વેઠ્યું}, \emph{teṇe ghaṇũ duḥkh veṭhyũ}.

\begin{cousinweb}
\textbf{\gu{વેઠવું}} (\emph{veṭhvũ}) is built on \emph{veṭh}, \textquotedblleft{}forced labour\textquotedblright{}, which CDIAL traces to Sanskrit \emph{viṣṭi} (CDIAL 11989); \textquotedblleft{}to endure hardship\textquotedblright{} is the verb's sense, which CDIAL does not give.
\end{cousinweb}

\begin{grammarlens}[title={What you've built}]
One more word for saying what happens if.
\end{grammarlens}

\subsection*{Your turn}
\begin{itemize}
\raggedright
  \item \emph{Say it:} \emph{veṭhvũ}
  \item \emph{Say it:} \emph{veṭhvũ}, once more
  \item \emph{Say it:} \emph{jo tame atyāre mahenat karśo to pachī duḥkh veṭhvũ nahī̃ paḍe.} — If you work hard now, you will not have to suffer later.
  \item \emph{Recall:} say the Gujarati for a misfortune, then the Gujarati for to spoil, then say \emph{veṭhvũ} again
\end{itemize}

\begin{tcolorbox}[breakable,colback=teal!4,colframe=teal!35!black,title={Before you move on}]
What does \gu{વેઠવું} mean? (\textquotedblleft{}To put up with (hardship)\textquotedblright{}.) Say it once more.
\end{tcolorbox}
\section[(dialogue)]{Looking back: from What the Figures Say to Danger and Harm}

\label{lesson:GU-R370-review-363-370}



\begin{quote}
Say the Gujarati for a figure and a calculation.
\end{quote}

\subsection*{Your turn}
Say these aloud:

\begin{itemize}
\raggedright
  \item a figure, a calculation, a finding, a count, to bring down
  \item a discovery, an observation, proven, a mystery, to uncover
  \item a principle, everyday practice, an inference, actual, to link
  \item a stipulation, possible, a likelihood, a conjecture, to want
  \item probably, without fail, definite, an assumption, to assess
  \item luck, destiny, a coincidence, to bear fruit, to turn out
  \item dread, terror, peace of mind, a misgiving, to frighten
  \item a peril, harm, a misfortune, to spoil, to put up with
\end{itemize}

\begin{tcolorbox}[breakable,colback=teal!4,colframe=teal!35!black,title={Before you move on}]
A figure? (\textbf{\gu{આંકડો}}, \emph{ā̃kḍo}.) A calculation? (\textbf{\gu{ગણતરી}}, \emph{gaṇtarī}.) A finding? (\textbf{\gu{તારણ}}, \emph{tāraṇ}.) A count? (\textbf{\gu{સંખ્યા}}, \emph{sankhyā}.) To bring down? (\textbf{\gu{ઘટાડવું}}, \emph{ghaṭāḍvũ}.) A discovery? (\textbf{\gu{શોધ}}, \emph{śodh}.) An observation? (\textbf{\gu{અવલોકન}}, \emph{avlokan}.) Proven? (\textbf{\gu{સિદ્ધ}}, \emph{siddh}.) A mystery? (\textbf{\gu{રહસ્ય}}, \emph{rahasya}.) To uncover? (\textbf{\gu{ઉઘાડવું}}, \emph{ughāḍvũ}.) A principle? (\textbf{\gu{સિદ્ધાંત}}, \emph{siddhā̃t}.) Everyday practice? (\textbf{\gu{વ્યવહાર}}, \emph{vyavhār}.) An inference? (\textbf{\gu{અનુમાન}}, \emph{anumān}.) Actual? (\textbf{\gu{વાસ્તવિક}}, \emph{vāstavik}.) To link? (\textbf{\gu{સાંકળવું}}, \emph{sā̃kaḷvũ}.) A stipulation? (\textbf{\gu{શરત}}, \emph{śarat}.) Possible? (\textbf{\gu{શક્ય}}, \emph{śakya}.) A likelihood? (\textbf{\gu{સંભાવના}}, \emph{sambhāvnā}.) A conjecture? (\textbf{\gu{અટકળ}}, \emph{aṭkaḷ}.) To want? (\textbf{\gu{ચાહવું}}, \emph{cāhvũ}.) Probably? (\textbf{\gu{સંભવતઃ}}, \emph{sambhavataḥ}.) Without fail? (\textbf{\gu{અચૂક}}, \emph{acūk}.) Definite? (\textbf{\gu{નિશ્ચિત}}, \emph{niścit}.) An assumption? (\textbf{\gu{ધારણા}}, \emph{dhārṇā}.) To assess? (\textbf{\gu{આંકવું}}, \emph{ā̃kvũ}.) Luck? (\textbf{\gu{નસીબ}}, \emph{nasīb}.) Destiny? (\textbf{\gu{ભાગ્ય}}, \emph{bhāgya}.) A coincidence? (\textbf{\gu{સંયોગ}}, \emph{sanyog}.) To bear fruit? (\textbf{\gu{ફળવું}}, \emph{phaḷvũ}.) To turn out? (\textbf{\gu{નીવડવું}}, \emph{nīvaḍvũ}.) Dread? (\textbf{\gu{ભય}}, \emph{bhay}.) Terror? (\textbf{\gu{દહેશત}}, \emph{daheśat}.) Peace of mind? (\textbf{\gu{નિરાંત}}, \emph{nirā̃t}.) A misgiving? (\textbf{\gu{આશંકા}}, \emph{āśankā}.) To frighten? (\textbf{\gu{ડરાવવું}}, \emph{ḍarāvvũ}.) A peril? (\textbf{\gu{ખતરો}}, \emph{khatro}.) Harm? (\textbf{\gu{હાનિ}}, \emph{hāni}.) A misfortune? (\textbf{\gu{આપત્તિ}}, \emph{āpatti}.) To spoil? (\textbf{\gu{બગાડવું}}, \emph{bagāḍvũ}.) To put up with? (\textbf{\gu{વેઠવું}}, \emph{veṭhvũ}.)
\end{tcolorbox}
